% ============================================
% DESEMPENHO DO DSB-SC
% Segue o diagrama do receptor: ponto 1, antes da detecção, com ruído
% passa-banda; ponto 2, depois da detecção, com tudo em banda base.
% ============================================

% Marcador numerado dos pontos de medida no diagrama
\tikzset{marca/.style={circle, inner sep=0pt, minimum size=3.6mm, font=\tiny\bfseries}}

% Diagrama compacto do receptor DSB-SC. Argumento: ponto destacado, 1 ou 2.
\newcommand{\diagDSB}[1]{%
\begin{tikzpicture}[font=\scriptsize,
  blk/.style={rectangle, draw, rounded corners, minimum height=0.6cm, minimum width=1.2cm, align=center},
  sm/.style={circle, draw, inner sep=0.5pt}]
\node (in) {$s(t)$};
\node[sm, right=0.5cm of in] (sum) {$+$};
\node[above=0.25cm of sum, cRed] (n) {$n(t)$};
\node[blk, right=0.5cm of sum] (bpf) {PF $2W$\\ em $f_c$};
\node[sm, right=1.3cm of bpf] (mix) {$\times$};
\node[above=0.25cm of mix] (lo) {$\cos(2\pi f_c t)$};
\node[blk, right=0.5cm of mix] (lpf) {PB $W$};
\node[right=1.3cm of lpf] (out) {$y(t)$};
\draw[->] (in) -- (sum);
\draw[->, cRed] (n) -- (sum);
\draw[->] (sum) -- (bpf);
\draw[->] (bpf) -- (mix);
\draw[->] (lo) -- (mix);
\draw[->] (mix) -- (lpf);
\draw[->] (lpf) -- (out);
\coordinate (p1) at ($(bpf.east)!0.5!(mix.west)$);
\coordinate (p2) at ($(lpf.east)!0.5!(out.west)$);
\ifnum#1=1
  \node[marca, fill=cBlue, text=white] at (p1) {1};
\else
  \node[marca, draw=gray!60, fill=white, text=gray] at (p1) {1};
\fi
\ifnum#1=2
  \node[marca, fill=cGreen, text=white] at (p2) {2};
\else
  \node[marca, draw=gray!60, fill=white, text=gray] at (p2) {2};
\fi
\end{tikzpicture}}

% Diagrama no topo dos frames de cálculo, com pouco espaço vertical
\newcommand{\topoDSB}[1]{{\centering\diagDSB{#1}\par}\vspace{2pt}}

\section{Modulação AM}

\subsection{DSB-SC}

\begin{frame}{Receptor DSB-SC}

\[
s(t) = A_c\, m(t)\cos(2\pi f_c t)
\]

\begin{center}
\begin{tikzpicture}[font=\small,
  block/.style={rectangle, draw, rounded corners, minimum height=0.9cm, minimum width=1.7cm, align=center},
  sum/.style={circle, draw, inner sep=1pt}]
\node (in) {$s(t)$};
\node[sum, right=0.7cm of in] (sum) {$+$};
\node[below=0.5cm of sum, cRed] (n) {$n(t)$};
\draw[->] (in) -- (sum);
\draw[->, cRed] (n) -- (sum);
\onslide<2->{
  \node[block, right=0.7cm of sum] (bpf) {Filtro PF\\ $2W$ em $f_c$};
  \draw[->] (sum) -- (bpf);
}
\onslide<3->{
  \node[sum, right=1.6cm of bpf] (mix) {$\times$};
  \node[below=0.5cm of mix] (lo) {$\cos(2\pi f_c t)$};
  \draw[->] (bpf) -- (mix);
  \draw[->] (lo) -- (mix);
}
\onslide<4->{
  \node[block, right=0.9cm of mix] (lpf) {Filtro PB\\ banda $W$};
  \node[right=1.6cm of lpf] (out) {$y(t)$};
  \draw[->] (mix) -- (lpf);
  \draw[->] (lpf) -- (out);
}
\onslide<5->{
  \node[nota=cBlue, above=0.4cm of bpf] (a) {deixa passar só a banda do DSB};
  \draw[setinha, cBlue] (a) -- (bpf.north);
  \node[nota=cGreen, right=0.4cm of lo] (b) {portadora local coerente};
  \draw[setinha, cGreen] (b.west) -- (lo.east);
}
\onslide<6->{
  \node[marca, fill=cBlue, text=white, minimum size=4.5mm, font=\scriptsize\bfseries] (m1) at ($(bpf.east)!0.5!(mix.west)$) {1};
  \node[marca, fill=cGreen, text=white, minimum size=4.5mm, font=\scriptsize\bfseries] (m2) at ($(lpf.east)!0.5!(out.west)$) {2};
}
\end{tikzpicture}
\end{center}

\vspace{0.2cm}
\only<1>{O sinal DSB-SC ocupa a faixa $f_c \pm W$: banda $2W$, o dobro da mensagem.}%
\only<2>{O filtro passa-faixa corta todo o ruído fora da banda do sinal.}%
\only<3>{A multiplicação pela portadora local traz o sinal de volta para a banda base.}%
\only<4>{O passa-baixas remove as componentes em $2f_c$.}%
\only<5>{Mesmas condições da banda base: potência recebida $S_i$ e ruído branco $N_0/2$.}%
\only<6->{Calculamos sinal, ruído e SNR em dois pontos. No ponto 1, antes da detecção, o ruído ainda é passa-banda. No ponto 2, depois da detecção, tudo está em banda base.}

\end{frame}

% ============================================
% PONTO 1: ANTES DA DETECÇÃO
% ============================================

\begin{frame}{Ponto 1: sinal antes da detecção}

\topoDSB{1}

\small\compacto
O passa-faixa deixa o DSB-SC passar intacto. A potência no ponto 1 é a potência recebida:
\onslide<2->{
\[
S_i = \overline{A_c^2\, m^2(t)\,\tikzmarknode{cq}{\cos^2(2\pi f_c t)}}
\only<3->{= \frac{A_c^2}{2}\,\overline{m^2} + \tikzmarknode{hf}{\frac{A_c^2}{2}\,\overline{m^2(t)\cos(4\pi f_c t)}}}
\]}%
\par\vspace{0.9cm}%
\anotacoes{
  \only<3->{
    \node[nota=cBlue, anchor=north east] (a) at ([xshift=0.4cm,yshift=-0.35cm]cq.south) {$\cos^2 x = \tfrac12 + \tfrac12\cos 2x$};
    \draw[setinha, cBlue] ([xshift=-0.4cm]a.north east) -- (cq.south);
  }
  \only<4->{
    \node[nota=cRed, anchor=north east] (b) at ([xshift=1.6cm,yshift=-0.35cm]hf.south) {média nula: $m^2$ varia devagar\\ perto de $\cos 4\pi f_c t$};
    \draw[setinha, cRed] ([xshift=-1.6cm]b.north east) -- (hf.south);
  }
}%
\onslide<5->{
\[
\boxed{\,S_i = \frac{A_c^2}{2}\,\overline{m^2}\,}
\]}

\end{frame}

% ============================================

\begin{frame}{Ponto 1: ruído antes da detecção}

\topoDSB{1}

\begin{columns}[T]
\column{0.52\textwidth}
\begin{tikzpicture}
\begin{axis}[onda, width=6cm, height=1.4cm, xmin=-4, xmax=4, ymin=0, ymax=1.3,
             xlabel={$f$}, ylabel={$S_n(f)$ no ponto 1}]
  \only<1>{\addplot[cRed, thick, domain=-4:4, samples=2] {0.6};}
  \only<2->{
    \addplot[cRed, thick, fill=cRed!20, samples=2] coordinates {(-3.1,0) (-3.1,0.6) (-1.9,0.6) (-1.9,0)};
    \addplot[cRed, thick, fill=cRed!20, samples=2] coordinates {(1.9,0) (1.9,0.6) (3.1,0.6) (3.1,0)};
    \node[font=\tiny, anchor=north] at (axis cs:2.5,0) {$f_c$};
    \node[font=\tiny, anchor=north] at (axis cs:-2.5,0) {$-f_c$};
  }
  \node[cRed, font=\tiny, anchor=south] at (axis cs:0,0.6) {$N_0/2$};
  \onslide<3->{
    \draw[cBlue, <->] (axis cs:1.9,0.85) -- node[above, font=\tiny] {$2W$} (axis cs:3.1,0.85);
    \draw[cBlue, <->] (axis cs:-3.1,0.85) -- node[above, font=\tiny] {$2W$} (axis cs:-1.9,0.85);
  }
\end{axis}
\end{tikzpicture}

\vspace{0.2cm}
\small
\only<1>{Ruído branco antes do filtro.}%
\only<2->{Ruído passa-banda: só sobram duas faixas, em torno de $\pm f_c$.}

\column{0.48\textwidth}
\small\compacto
\onslide<3->{Potência: área das duas faixas.}
\onslide<4->{
\[
\overline{n^2(t)} = \tikzmarknode{dois}{2}\times \frac{N_0}{2}\times \tikzmarknode{larg}{2W}
\]}%
\par\vspace{0.7cm}%
\anotacoes{\only<5->{
  \node[nota=cRed, anchor=north] (a) at ([xshift=-0.5cm,yshift=-0.4cm]dois.south) {duas faixas};
  \draw[setinha, cRed] (a.north) -- (dois.south);
  \node[nota=cBlue, anchor=north] (b) at ([xshift=0.9cm,yshift=-0.4cm]larg.south) {altura $\times$ largura};
  \draw[setinha, cBlue] (b.north) -- (larg.south);
}}%
\onslide<6->{
\[
\boxed{\,\overline{n^2(t)} = 2N_0W\,}
\]}%
\onslide<7->{O dobro do ruído da banda base.}
\end{columns}

\end{frame}

% ============================================

\begin{frame}{Ponto 1: SNR antes da detecção}

\topoDSB{1}

\small\compacto
Sinal e ruído no mesmo ponto, os dois ainda em torno de $f_c$:
\onslide<2->{
\[
\left(\frac{S}{N}\right)_1 = \frac{S_i}{\overline{n^2}}
\only<3->{= \frac{S_i}{2N_0W}}
\only<4->{= \frac12\,\tikzmarknode{gm}{\frac{S_i}{N_0W}}}
\only<5->{= \frac{\gamma}{2}}
\]}%
\par\vspace{0.9cm}%
\anotacoes{\only<4->{
  \node[nota=cGreen, anchor=north] (a) at ([yshift=-0.35cm]gm.south) {$\gamma$: SNR da banda base com a mesma $S_i$};
  \draw[setinha, cGreen] (a.north) -- (gm.south);
}}%
\onslide<6->{Antes de detectar, o DSB-SC está 3\,dB abaixo da banda base: entrou o dobro de ruído.}

\onslide<7->{Agora o sinal atravessa o detector. Para isso precisamos escrever o ruído passa-banda de um jeito que a multiplicação por $\cos$ saiba tratar.}

\end{frame}

% ============================================
% PONTO 2: DEPOIS DA DETECÇÃO
% ============================================

\begin{frame}{Ruído passa-banda em fase e quadratura}

\topoDSB{1}

\small\compacto
O ruído do ponto 1 com duas componentes de banda base, como $s_I$ e $s_Q$ no AM:
\onslide<2->{
\[
n(t) = \tikzmarknode{nc}{n_c(t)}\cos(2\pi f_c t) - \tikzmarknode{ns}{n_s(t)}\sin(2\pi f_c t)
\]}%
\par\vspace{0.7cm}%
\anotacoes{
  \only<3->{
    \node[nota=cBlue, anchor=north east] (a) at ([xshift=-0.3cm,yshift=-0.3cm]nc.south) {em fase, banda $W$};
    \draw[setinha, cBlue] (a.north) -- (nc.south);
  }
  \only<4->{
    \node[nota=cOrange, anchor=north west] (b) at ([xshift=0.3cm,yshift=-0.3cm]ns.south) {em quadratura, banda $W$};
    \draw[setinha, cOrange] (b.north) -- (ns.south);
  }
}%
\onslide<5->{Potência, com $\overline{\cos^2} = \overline{\sin^2} = \tfrac12$ e produto cruzado de média nula:}
\onslide<6->{
\[
\overline{n^2} = \tfrac12\,\overline{n_c^2} + \tfrac12\,\overline{n_s^2}
\only<7->{\quad\Rightarrow\quad \tikzmarknode{res}{\overline{n_c^2} = \overline{n_s^2} = \overline{n^2} = 2N_0W}}
\]}%
\par\vspace{0.2cm}%
\anotacoes{\only<7->{
  \node[nota=cRed, anchor=north] (c) at ([yshift=-0.3cm]res.south) {$n_c$ e $n_s$ dividem a potência por igual};
  \draw[setinha, cRed] (c.north) -- (res.south);
}}%

\end{frame}

% ============================================

\begin{frame}{Ponto 2: atravessando o detector}

\topoDSB{2}

\small\compacto
\onslide<1->{Entrada do detector, agrupando o que multiplica cada portadora:}
\[
r(t) = \big[A_c m + n_c\big]\cos(2\pi f_c t) - n_s\sin(2\pi f_c t)
\]
\onslide<2->{Multiplicando por $\cos(2\pi f_c t)$:}
\onslide<2->{
\[
r\cos(2\pi f_c t) = \big[A_c m + n_c\big]\tikzmarknode{c2}{\cos^2(2\pi f_c t)} - n_s\,\tikzmarknode{sc}{\sin(2\pi f_c t)\cos(2\pi f_c t)}
\]}%
\par\vspace{0.6cm}%
\anotacoes{
  \only<3->{
    \node[nota=cBlue, anchor=north] (a) at ([yshift=-0.25cm]c2.south) {$\tfrac12 + \tfrac12\cos 4\pi f_c t$};
    \draw[setinha, cBlue] (a.north) -- (c2.south);
  }
  \only<4->{
    \node[nota=cOrange, anchor=north] (b) at ([yshift=-0.25cm]sc.south) {$\tfrac12\sin 4\pi f_c t$: só em $2f_c$};
    \draw[setinha, cOrange] (b.north) -- (sc.south);
  }
}%
\onslide<5->{Depois do passa-baixas: \quad $\boxed{\,y(t) = \tfrac12 A_c\, m(t) + \tfrac12 n_c(t)\,}$}

\onslide<6->{O $n_s$ sumiu: o detector coerente rejeita o ruído em quadratura.}

\end{frame}

% ============================================

\begin{frame}{Ponto 2: sinal depois da detecção}

\topoDSB{2}

\small\compacto
\[
y(t) = \tikzmarknode{ys}{\tfrac12 A_c\, m(t)} + \tfrac12 n_c(t)
\]
\par\vspace{0.5cm}%
\anotacoes{
  \node[nota=cBlue, anchor=north] (a) at ([yshift=-0.25cm]ys.south) {parcela de sinal};
  \draw[setinha, cBlue] (a.north) -- (ys.south);
}%
\onslide<2->{O fator $\tfrac12$ entra ao quadrado na potência:}
\onslide<2->{
\[
S_o = \frac{A_c^2}{4}\,\overline{m^2}
\only<3->{= \frac12\cdot\tikzmarknode{si}{\frac{A_c^2}{2}\,\overline{m^2}}}
\only<4->{= \frac{S_i}{2}}
\]}%
\par\vspace{0.2cm}%
\anotacoes{\only<3->{
  \node[nota=cBlue, anchor=west] (b) at ([xshift=1.6cm]si.east) {$S_i$ do ponto 1};
  \draw[setinha, cBlue] (b.west) -- (si.east);
}}%
\onslide<5->{Mensagem em banda base, com metade da potência que tinha no ponto 1.}

\end{frame}

% ============================================

\begin{frame}{Ponto 2: PSD do ruído $n_c$}

\begin{columns}[T]
\column{0.55\textwidth}
\begin{tikzpicture}
\begin{axis}[onda, width=6.2cm, height=0.9cm, xmin=-5.8, xmax=5.8, ymin=0, ymax=1.3,
             xlabel={$f$}, ylabel={$S_n(f-f_c)$}]
  \addplot[cRed, thick, fill=cRed!20, samples=2] coordinates {(-0.6,0) (-0.6,0.5) (0.6,0.5) (0.6,0)};
  \addplot[cRed, thick, fill=cRed!20, samples=2] coordinates {(4.4,0) (4.4,0.5) (5.6,0.5) (5.6,0)};
  \node[font=\tiny, anchor=north] at (axis cs:5,0) {$2f_c$};
\end{axis}
\onslide<2->{
\begin{axis}[onda, width=6.2cm, height=0.9cm, at={(0,-1.8cm)}, xmin=-5.8, xmax=5.8, ymin=0, ymax=1.3,
             xlabel={$f$}, ylabel={$S_n(f+f_c)$}]
  \addplot[cRed, thick, fill=cRed!20, samples=2] coordinates {(-0.6,0) (-0.6,0.5) (0.6,0.5) (0.6,0)};
  \addplot[cRed, thick, fill=cRed!20, samples=2] coordinates {(-5.6,0) (-5.6,0.5) (-4.4,0.5) (-4.4,0)};
  \node[font=\tiny, anchor=north] at (axis cs:-5,0) {$-2f_c$};
\end{axis}
}
\onslide<3->{
\begin{axis}[onda, width=6.2cm, height=0.9cm, at={(0,-3.6cm)}, xmin=-5.8, xmax=5.8, ymin=0, ymax=1.3,
             xlabel={$f$}, ylabel={$S_{n_c}(f)$}]
  \addplot[cPurple, thick, fill=cPurple!20, samples=2] coordinates {(-0.6,0) (-0.6,1.0) (0.6,1.0) (0.6,0)};
  \node[font=\tiny, anchor=north] at (axis cs:-0.6,0) {$-W$};
  \node[font=\tiny, anchor=north] at (axis cs:0.6,0) {$W$};
  \node[cPurple, font=\tiny, anchor=west] at (axis cs:0.7,1.0) {$N_0$};
\end{axis}
}
\end{tikzpicture}

\column{0.45\textwidth}
\small\compacto
O detector desloca o ruído do ponto 1 para a banda base.
\vspace{0.2cm}

\only<1>{Multiplicar por $\cos$ desloca a PSD de $+f_c$: uma faixa cai em torno de zero, a outra em $2f_c$.}%
\only<2>{E de $-f_c$: outra faixa cai em torno de zero.}%
\only<3->{O passa-baixas guarda só $|f|\le W$, onde as duas faixas se somam:
\[
S_{n_c}(f) = \frac{N_0}{2}+\frac{N_0}{2} = N_0, \quad |f|\le W
\]}%
\onslide<4->{Área: $N_0\cdot 2W = 2N_0W$, a mesma potência do slide anterior.}
\end{columns}

\end{frame}

% ============================================

\begin{frame}{Ponto 2: ruído depois da detecção}

\topoDSB{2}

\begin{columns}[T]
\column{0.55\textwidth}
\small\compacto
\[
y(t) = \tfrac12 A_c\, m(t) + \tikzmarknode{yn}{\tfrac12 n_c(t)}
\]
\par\vspace{0.4cm}%
\anotacoes{
  \node[nota=cRed, anchor=west] (a) at ([xshift=0.5cm]yn.east) {parcela de ruído};
  \draw[setinha, cRed] (a.west) -- (yn.east);
}%
\onslide<2->{O fator $\tfrac12$ divide a PSD por 4:}
\onslide<2->{
\[
S_{n_o}(f) = \tfrac14\, S_{n_c}(f)
\only<3->{= \frac{N_0}{4}, \quad |f|\le W}
\]}%
\onslide<4->{
\[
N_o = \frac{N_0}{4}\cdot 2W
\only<5->{= \frac{N_0W}{2}}
\]}%
\onslide<6->{Ruído de banda base, com metade de $N_0W$.}

\column{0.45\textwidth}
\onslide<3->{
\begin{tikzpicture}
\begin{axis}[onda, width=4.4cm, height=1.6cm, xmin=-2.2, xmax=2.2, ymin=0, ymax=1.3,
             xlabel={$f$}, ylabel={$S_{n_o}(f)$}]
  \addplot[cRed, thick, fill=cRed!20, samples=2] coordinates {(-1,0) (-1,0.5) (1,0.5) (1,0)};
  \node[font=\tiny, anchor=north] at (axis cs:-1,0) {$-W$};
  \node[font=\tiny, anchor=north] at (axis cs:1,0) {$W$};
  \node[cRed, font=\tiny, anchor=east] at (axis cs:-1,0.5) {$N_0/4$};
\end{axis}
\end{tikzpicture}}
\end{columns}

\end{frame}

% ============================================

\begin{frame}{Ponto 2: SNR depois da detecção}

\topoDSB{2}

\small\compacto
Sinal e ruído no ponto 2, os dois em banda base, com \textcolor{cBlue}{$S_o = S_i/2$} e \textcolor{cRed}{$N_o = N_0W/2$}:
\onslide<2->{
\[
\left(\frac{S}{N}\right)_2 = \frac{\color{cBlue}S_o}{\color{cRed}N_o}
\only<3->{= \frac{\color{cBlue}S_i/2}{\color{cRed}N_0W/2}}
\only<4->{= \frac{S_i}{N_0W}}
\]}%
\onslide<5->{Resultado: \quad $\boxed{\,\left(\frac{S}{N}\right)_{2} = \gamma\,}$}

\onslide<6->{O DSB-SC com detecção coerente tem a mesma SNR de saída da banda base.}

\end{frame}

% ============================================

\begin{frame}{Comparando os dois pontos}

\begin{center}
\begin{tikzpicture}[font=\small,
  block/.style={rectangle, draw, rounded corners, minimum height=0.9cm, minimum width=1.7cm, align=center},
  sum/.style={circle, draw, inner sep=1pt},
  medida/.style={rounded corners=3pt, inner sep=4pt, font=\footnotesize, align=left, line width=0.7pt}]
\node (in) {$s(t)$};
\node[sum, right=0.7cm of in] (sum) {$+$};
\node[above=0.5cm of sum, cRed] (n) {$n(t)$};
\node[block, right=0.7cm of sum] (bpf) {Filtro PF\\ $2W$ em $f_c$};
\node[sum, right=1.6cm of bpf] (mix) {$\times$};
\node[above=0.5cm of mix] (lo) {$\cos(2\pi f_c t)$};
\node[block, right=0.8cm of mix] (lpf) {Filtro PB\\ banda $W$};
\node[right=1.6cm of lpf] (out) {$y(t)$};
\draw[->] (in) -- (sum);
\draw[->, cRed] (n) -- (sum);
\draw[->] (sum) -- (bpf);
\draw[->] (bpf) -- (mix);
\draw[->] (lo) -- (mix);
\draw[->] (mix) -- (lpf);
\draw[->] (lpf) -- (out);
\coordinate (p1) at ($(bpf.east)!0.5!(mix.west)$);
\coordinate (p2) at ($(lpf.east)!0.5!(out.west)$);
\node[marca, fill=cBlue, text=white, minimum size=4.5mm, font=\scriptsize\bfseries] at (p1) {1};
\node[marca, fill=cGreen, text=white, minimum size=4.5mm, font=\scriptsize\bfseries] at (p2) {2};
\onslide<2->{
  \draw[cBlue, dashed] ($(p1)+(0,-0.25)$) -- ++(0,-0.75);
  \node[medida, draw=cBlue, text=cBlue, fill=cBlue!6, anchor=north] (m1) at ($(p1)+(0,-1.0)$) {%
    passa-banda\\[2pt]
    sinal: $S_i$\\[2pt]
    ruído: $2N_0W$\\[2pt]
    SNR: $\gamma/2$};
}
\onslide<3->{
  \draw[cGreen, dashed] ($(p2)+(0,-0.25)$) -- ++(0,-0.75);
  \node[medida, draw=cGreen, text=cGreen, fill=cGreen!6, anchor=north] (m2) at ($(p2)+(0,-1.0)$) {%
    banda base\\[2pt]
    sinal: $S_i/2$\\[2pt]
    ruído: $N_0W/2$\\[2pt]
    SNR: $\gamma$};
}
\onslide<4->{
  \draw[setinha, cOrange] (m1.east) -- node[above, font=\scriptsize, cOrange] {$\times 2$} (m2.west);
}
\end{tikzpicture}
\end{center}

\vspace{0.1cm}
\small
\only<1>{Os dois pontos de medida no mesmo receptor.}%
\only<2>{No ponto 1, o ruído ocupa a banda $2W$ do DSB.}%
\only<3>{No ponto 2, o sinal cai pela metade e o ruído cai a um quarto.}%
\only<4->{O detector dobra a SNR: ganho de detecção 2. A comparação com a banda base é feita no ponto 2.}

\end{frame}

% ============================================

\begin{frame}{Nem ganha nem perde: por quê?}

\small
\begin{itemize}
\item<1-> O detector coerente aproveita todo o sinal, que está em fase com a portadora local.
\item<2-> Do ruído passa-banda, aproveita só $n_c$ e descarta $n_s$, que está em quadratura.
\item<3-> Esse ganho de 2 só compensa o dobro de ruído que entrou pela banda $2W$.
\end{itemize}

\onslide<4->{
\begin{center}
\begin{tabular}{lcccc}
\toprule
 & banda & ruído no ponto 1 & SNR no ponto 1 & SNR de saída \\
\midrule
banda base & $W$ & $N_0W$ & $\gamma$ & $\gamma$ \\
DSB-SC & $2W$ & $2N_0W$ & $\gamma/2$ & $\gamma$ \\
\bottomrule
\end{tabular}
\end{center}}

\onslide<5->{O DSB-SC paga com o dobro de banda e entrega a mesma qualidade da banda base.}

\end{frame}

% ============================================

\begin{frame}{Erro de fase da portadora local}

\small\compacto
Se o oscilador local é $\cos(2\pi f_c t + \phi)$, a mesma conta dá:
\onslide<2->{
\[
y(t) = \tfrac12\big[A_c m + n_c\big]\tikzmarknode{cf}{\cos\phi} + \tfrac12\, n_s \tikzmarknode{sf}{\sin\phi}
\]}%
\par\vspace{0.6cm}%
\anotacoes{
  \only<3->{
    \node[nota=cBlue, anchor=north east] (a) at ([xshift=-0.2cm,yshift=-0.3cm]cf.south) {sinal atenuado};
    \draw[setinha, cBlue] (a.north) -- (cf.south);
  }
  \only<4->{
    \node[nota=cOrange, anchor=north west] (b) at ([xshift=0.2cm,yshift=-0.3cm]sf.south) {$n_s$ volta};
    \draw[setinha, cOrange] (b.north) -- (sf.south);
  }
}%
\onslide<5->{
\[
S_o = \frac{A_c^2\overline{m^2}}{4}\cos^2\phi, \quad N_o = \frac14\big(2N_0W\cos^2\phi + 2N_0W\sin^2\phi\big)
\only<6->{= \frac{N_0W}{2}}
\]}%
\onslide<7->{
\[
\left(\frac{S}{N}\right)_o = \gamma\cos^2\phi
\]}%
\onslide<8->{O ruído total não muda, só o sinal encolhe. Com $\phi = 90^\circ$, a saída é só ruído.}

\end{frame}

% ============================================

\begin{frame}{Exemplo: enlace DSB-SC}

Um transmissor DSB-SC entrega $P_T = 2$\,W a um enlace com perda de 90\,dB. A mensagem tem banda $W = 5$\,kHz e o receptor tem $N_0 = 2\times10^{-17}$\,W/Hz.

\begin{enumerate}
\item<2-> Qual a SNR antes da detecção, no ponto 1?
\item<3-> Qual a SNR depois da detecção coerente, no ponto 2?
\item<4-> E se a portadora local tiver erro de fase de $30^\circ$?
\end{enumerate}

\vspace{0.4cm}
\onslide<2->{
\begin{center}
\begin{tikzpicture}[font=\small,
  block/.style={rectangle, draw, rounded corners, minimum height=0.8cm, minimum width=1.5cm, align=center},
  sum/.style={circle, draw, inner sep=1pt}]
\node[block] (tx) {Tx DSB\\\scriptsize $2$\,W};
\node[block, right=0.7cm of tx] (ch) {Enlace\\\scriptsize $-90$\,dB};
\node[sum, right=0.7cm of ch] (sum) {$+$};
\node[below=0.4cm of sum, cRed, font=\scriptsize] (n) {$N_0/2$};
\node[block, right=0.7cm of sum] (bpf) {PF\\\scriptsize $10$\,kHz};
\node[block, right=0.7cm of bpf] (det) {Detector\\\scriptsize coerente};
\node[right=0.5cm of det] (out) {$y(t)$};
\draw[->] (tx) -- (ch);
\draw[->] (ch) -- node[above, font=\scriptsize] {$S_i$} (sum);
\draw[->, cRed] (n) -- (sum);
\draw[->] (sum) -- (bpf);
\draw[->] (bpf) -- (det);
\draw[->] (det) -- (out);
\end{tikzpicture}
\end{center}}

\end{frame}

% ============================================

\begin{frame}{Resolução}

\small\compacto
\onslide<1->{Potência recebida, com 90\,dB $= 10^9$:}
\[
S_i = \frac{2}{10^9} = 2\times10^{-9}\ \text{W}
\]
\onslide<2->{SNR de referência:}
\onslide<2->{
\[
\gamma = \frac{S_i}{N_0W}
\only<3->{= \frac{2\times10^{-9}}{2\times10^{-17}\cdot 5\times10^{3}}}
\only<4->{= 2\times10^{4} \approx 43\ \text{dB}}
\]}%
\onslide<5->{No ponto 1, antes da detecção: $\gamma/2 = 10^4 = 40$\,dB.}

\onslide<6->{No ponto 2, depois da detecção: $\gamma \approx 43$\,dB, 3\,dB de ganho de detecção.}

\onslide<7->{Com $\phi = 30^\circ$:}
\onslide<7->{
\[
\gamma\cos^2 30^\circ
\only<8->{= 2\times10^4\cdot 0{,}75}
\only<9->{= \tikzmarknode{res}{1{,}5\times10^4 \approx 41{,}8\ \text{dB}}}
\]}%
\par\vspace{0.6cm}%
\anotacoes{\only<10->{
  \node[nota=cRed, anchor=west] (a) at ([xshift=0.8cm]res.east) {perde $1{,}25$\,dB};
  \draw[setinha, cRed] (a.west) -- (res.east);
}}%

\end{frame}

% ============================================
